\section{Introduction}
\label{sec:intro}

With the rapid development of e-commerce, on-demand local delivery, and smart logistics, urban delivery systems are facing stricter timeliness requirements and more complex operational constraints. Compared with traditional fuel vehicles, new energy logistics vehicles offer advantages such as lower emissions, higher energy efficiency, and policy support, and are becoming an important part of urban delivery. However, they also introduce new challenges, including limited battery capacity, time-consuming charging, and limited charging-station resources with possible congestion, which make vehicle routing and task scheduling more difficult.

In real logistics scenarios, delivery tasks are usually not fully known in advance but arrive dynamically over time. Under constraints such as limited fleet size, vehicle load limits, continuous battery consumption, task deadlines, and charging-station queues, managers must decide in real time which vehicle should serve which task, how to plan routes, when to charge, and whether multiple vehicles should collaborate on the same task. Therefore, this is not simply a shortest-path problem, but a comprehensive problem involving graph search, resource-constrained scheduling, dynamic decision-making, and system simulation. Based on this background, we aim to design and implement a complete system platform to simulate and analyze the performance of different scheduling strategies in new energy logistics scenarios.
\subsection{Dynamic Modeling of Vehicles and Tasks}

In this project, vehicles, tasks, depots, and charging stations are modeled in a unified way at the entity level. Each vehicle includes attributes such as current location, battery level, battery capacity, maximum load capacity, current load, speed, and energy consumption per unit distance. Each task includes information such as release time, deadline, task location, cargo weight, and status. During the simulation, tasks are released dynamically as time progresses, rather than being provided to the scheduler all at once, so that the system can better reflect the online decision-making process in real delivery scenarios.

\subsection{Graph-Based Road Network and Path Planning}

This project uses graphs as the core data structure and models the road network as a weighted graph, where nodes represent road intersections, depots, task points, and charging stations, and edges represent road connections together with their distance or travel cost. For path planning, the system implements Dijkstra’s algorithm as a stable and reliable shortest-path baseline, while A* is introduced to improve goal-directed search efficiency. RRT is also included as an extended experimental algorithm to demonstrate the scalability of the system in path planning. With these algorithms, the system can support shortest-path queries, reachability checking, return-to-depot feasibility checking, and nearest charging station search.

\subsection{Reward and Charging Constraints}
This project implements a task reward calculation and overdue penalty mechanism in the simulation environment, and models each charging station as a resource node with a fixed number of chargers, occupancy status, and a waiting queue. The system not only checks whether a vehicle can complete a task and return to the depot, but also automatically triggers charging-station selection when the battery level is low. In station selection, it considers not only distance, but also queue length and station occupancy, so as to better reflect the practical complexity of new energy logistics delivery.
%-------------------------------------------------------------------------
\subsection{Scheduling, Experiments, and Visualization}
This project implements several heuristic scheduling strategies, including nearest-task-first, maximum-weight-first, and earliest-deadline-first, and supports different scenario settings such as small-scale, medium-scale, and large-scale cases. In addition, the project further explores the use of Q-learning for dynamic scheduling decisions and builds an offline MILP module to solve global plans for small-scale scenarios. Based on these components, the project also provides real-time visualization of the map, tasks, vehicles, charging stations, and statistical indicators through a web frontend, so that the system can be used not only for simulation but also for demonstration and result analysis.
\subsection{Summary}
This project adopts an overall architecture of “simulation engine + algorithm modules + real-time interface + web-based visualization frontend.” At the simulation level, the system maintains the graph-based road network, vehicle states, task states, charging station states, and the logic of time advancement. At each time step, it sequentially performs task release, task expiration checking, scheduling decisions, vehicle movement, battery consumption, charging queue management, and charging completion, forming a complete discrete-time simulation loop. At the algorithm level, path planning and task scheduling are designed as separate modules: path planning handles shortest-path search, target reachability, and energy feasibility analysis, while scheduling strategies determine dispatch actions based on visible tasks, vehicle states, and system constraints. This modular design allows different algorithms to be compared fairly in the same environment and makes future extensions easier.

To address the course requirement for advanced methods, the project further implements two extended approaches. One is a Q-learning-based hyper-heuristic strategy selection method, which learns which low-level heuristic rule should be used in different situations through discrete state encoding. The other is an MILP-based offline solver, which computes globally near-optimal routes and scheduling plans for small-scale cases under full-information settings, and compares them with the performance of online strategies. In addition, the system provides real-time visualization through frontend-backend interaction: the backend updates the environment and records logs, while the frontend displays map nodes, vehicle movements, task-state changes, charging station queues, and statistical indicators. Users can also switch scheduling strategies, replay offline solutions, and perform Q-learning training and inference through the interface, which improves both interactivity and presentation quality.
